%% fig_comm.tex — figure definitions, one \newcommand per figure.
%% Call the command (e.g. \figPlaceholder) at the point in the section where
%% the float should be anchored. This keeps section files readable and makes
%% figures easy to reorder or comment out.

\newcommand{\figPlaceholder}{%
  \begin{figure}[t]
    \centering
    \begin{tikzpicture}
      \draw[kaistblue, very thick, rounded corners]
        (0,0) rectangle (10,3.2);
      \node[align=center, text=kaistbluedark, font=\sffamily\large]
        at (5,1.6) {Figure placeholder\\[2pt]
        \footnotesize replace with \texttt{\textbackslash includegraphics\{figs/...\}}};
    \end{tikzpicture}
    \caption{A placeholder figure drawn with Ti\emph{k}Z so the template compiles
    with no external image files. Swap in your own graphic from \texttt{figs/}.}
    \label{fig:placeholder}
  \end{figure}%
}
